\begin{abstract}
Diffusion models on SE(3) handle the pose ambiguity of symmetric and occluded objects by sampling. Each denoising step, however, is a
feed-forward function of image features and the noisy pose, so the network never checks whether its current hypothesis explains the
image. Render-and-compare refiners do perform this check. They need a renderer and a textured 3D model, though, and run as a separate
stage after an initial estimate. We propose \method, a pose-estimation architecture that carries out this comparison in a learned feature
space inside every denoising step. Its world model is a feature field of the object: tri-planes and a small MLP, fitted once to the
frozen DINOv2 tokens of posed training crops. It renders by volume rendering the tokens that the object would produce at any pose
hypothesis, in the observation's crop. The field needs no reference image, CAD model or RGB reconstruction, and its 3D structure is
learned only from the tokens. At every denoising step and for every sample, the denoiser receives the observed tokens, the tokens
rendered at its current hypothesis, their difference, and the field's latent at the surface each ray hits. A control that removes only
the rendered quantities isolates the effect of the comparison. We evaluate on T-LESS under the BOP protocol, on the SPEED+ spacecraft
benchmark with its official score, and against LiePose trained with the same budget. A retrieval-referenced world model, which
predicts tokens from the nearest training view, is studied in the supplement.
\TODO{one sentence of results, written only after the tables are filled.}
\end{abstract}
